\documentclass[10pt,a4paper]{amsart}
\usepackage[margin=19mm]{geometry}
\usepackage{amsmath,amssymb,mathtools}
\usepackage[scr=rsfs,cal=euler]{mathalfa}
\usepackage{xfrac}
\usepackage[unicode,hidelinks,bookmarks=false]{hyperref}
\newcommand{\ie}{i.e.\ }
\newcommand{\cf}{cf.\ }
\newcommand{\resp}{respectively\ }
\newcommand{\Subsection}{\subsection}
\providecommand{\nobreakdash}{\nobreak\hbox{-}}
\setlength{\emergencystretch}{2em}
\multlinegap0pt
\usepackage{amssymb}
\usepackage{xcolor}
\usepackage{tikz}
\usepackage{enumerate}
\newtheorem{theorem}{Theorem}
\newcommand{\ad}{\textrm{ad}}
\title{Parabolic Lie algebroid connections on parabolic principal bundles over curves}
\author{Indranil Biswas, Pritthijit Biswas}
\date{Source: arXiv:2609.01402v1 (CC BY 4.0)}
\begin{document}
\maketitle

\noindent\emph{Source and licence.} Parabolic Lie algebroid connections on parabolic principal bundles over curves. Version arXiv:2609.01402v1 (CC BY 4.0), distributed by arXiv under the Creative Commons Attribution 4.0 International licence, http://creativecommons.org/licenses/by/4.0/, as recorded in the arXiv licence metadata for this exact version and shown on the abstract page https://arxiv.org/abs/2609.01402v1. Full licence notice: \texttt{licenses/arxiv-2609.01402-CC-BY-4.0.txt}.\par\smallskip

\noindent\textit{Editorial scope.} Counted unit: the article's theorem MAIN (source0.tex lines 1218--1227). The article's complete original proof of it is delivered with it (source0.tex lines 1229--1254, one \texttt{proof} environment of 26 lines). No mathematics is written, repaired or re-derived: the statement and the proof are the article's own lines, in the article's own notation, and no retained line is substituted or re-flowed.  Retained uncounted as the source-local context the proof needs: lines 607--629; lines 826--834; lines 1208--1210.  Nothing was omitted from any retained span, so the package carries no omission record.  No retained line contains a citation, so the wrapper carries no bibliography and no bibliography entry.  No retained line contains a figure environment, an image-inclusion command or drawing code, so no image is delivered and the package declares no asset dependency.  Editorial formatting only, in addition: an AMS \texttt{amsart} preamble that restates the article's own declarations and macros used by the retained lines, namely the environments \texttt{theorem} on separate counters, together with the standard packages those lines need.  The retained environments print in this document as Theorem 1 in order of appearance, whereas the article prints the same items with the numbering of its own full document.  The complete native source of the article as distributed in the arXiv e-print for this exact version is bundled unchanged as \texttt{sources/209108/source0.tex}.
\begin{theorem}\label{char}\mbox{}
\begin{enumerate}
\item There are natural equivalences of categories $\Phi\,:\,\mathscr{C}_{\Gamma}\,\overset{\sim}\longrightarrow
\, \mathscr{C}$ and $\Psi\,:\,\mathscr{V}_{\Gamma}\,\overset{\sim}\longrightarrow \,\mathscr{V}$. 

\item The following equalities of functors hold:
$$\Psi\circ\mathbb{A}\ = \ \mathbb{A}_{*}\circ\Phi\ \ \text{ and }\ \ \Psi\circ\mathrm{ad}\ =\ \mathrm{ad}_{*}\circ\Phi .$$

\item There is a natural isomorphism $\zeta\,:\,TX_{*}\, \overset{\sim}\longrightarrow\, \Psi(TY)$
given by the differential of $\sigma$. Moreover, for any $E_{G}\,\in\, \mathscr{C}_{\Gamma}$,
$\Psi$ takes the equivariant Atiyah exact sequence for 
$E_{G}$ to the parabolic Atiyah exact sequence for $\Phi(E_{G})\,\in\,\mathscr{C}$.

\item There is an equivalence of categories between the parabolic Lie algebroids on $(X,\,S)$ and the
$\Gamma$--equivariant Lie algebroids on $Y$. 

\item Take $E_{G}\,\in\, \mathscr{C}_{\Gamma}$ and a $\Gamma$--equivariant Lie algebroid
$(V,\,\phi)$ on $Y$. Giving a $\Gamma$--equivariant holomorphic $(V,\,\phi)$--connection on $E_{G}$ is equivalent
to giving a parabolic $(V_{*},\, \phi_{*})$--connection on $\mathcal{E}_{G}\,:=\,\Phi(E_{G})\,\in
\, \mathscr{C}$, where $(V_{*},\, \phi_{*})$ is the parabolic Lie algebroid on $X$ corresponding to
$(V,\, \phi)$.
\end{enumerate}
\end{theorem}

\begin{theorem}\label{4.1}
Take a compact connected Riemann surface $X$ and also a connected complex reductive affine algebraic group
$G$. Let $E_{G}$ be an equivariant holomorphic principal $G$--bundle on $X$ equipped
with an equivariant holomorphic reduction of structure group of $E_{P}\,\subset\, E_{G}$, to some parabolic
subgroup $P\, \subset\, G$, that is equivariantly infinitesimally rigid. Take any equivariant holomorphic
Lie algebroid $(V,\,\phi)$ on $X$ that satisfies the condition that the anchor map $\phi$ is not
surjective. Then
the holomorphic principal $P$--bundle $E_{P}$ admits an equivariant holomorphic $(V,\,\phi)$--connection.
\end{theorem}

\begin{equation}\label{coh}
H^{k}(Y,\ \ad(E_{G})/\ad(E_{P}))^{\Gamma}\,\ \cong\,\ H^{k}(X,\
\ad_{*}(\mathcal{E}_{G})/\ad_{*}(\mathcal{E}_{P})).\end{equation}

\begin{theorem}\label{MAIN}
Take a compact connected Riemann surface $X$ with a finite set of marked points $S\, \subset\, X$.
Take a connected complex reductive affine algebraic group $G$ and a parabolic subgroup $P\, \subset\, G$.
Let $\mathcal{E}_{G}$ be a parabolic principal $G$--bundle on $(X,\,S)$ equipped with a
parabolic reduction of structure group $\mathcal{E}_{P}\, \subset\, \mathcal{E}_{G}$ of
$\mathcal{E}_{G}$ to $P$ such that the reduction is parabolically infinitesimally rigid.
Then for any parabolic Lie algebroid $(V_{*},\, \phi_{*})$ on $(X,\,S)$ such that
$\phi_{*}$ is not surjective, there is a holomorphic parabolic $(V_{*},\, \phi_{*})$--connection
on $\mathcal{E}_{P}$.
\end{theorem}

\begin{proof}
Recall Theorem \ref{char}. There is a ramified Galois covering map $\sigma\,:\, Y\,\longrightarrow\, X$,
whose branch locus is $S$, such that $\mathcal{E}_{G}$ corresponds to a
$\Gamma$--equivariant holomorphic principal $G$--bundle $E_G$ on $Y$, where $\Gamma\,=\, \text{Gal}(\sigma)$
is the Galois group for $\sigma$. Let $E_{P}\,\subset\, E_{G}$ be the $\Gamma$--equivariant holomorphic
reduction of structure group of $E_{G}$ to $P$ given by reduction $\mathcal{E}_{P}\, \subset\, \mathcal{E}_{G}$ of
$\mathcal{E}_{G}$ to $P$. Since the reduction $\mathcal{E}_{P}\, \subset\, \mathcal{E}_{G}$
is parabolically infinitesimally rigid, from \eqref{coh} we conclude that the reduction
$E_{P}\,\subset\, E_{G}$ is equivariantly infinitesimally rigid.

Let $(V,\,\phi)$ be the $\Gamma$--equivariant holomorphic Lie algebroid on $Y$ that corresponds to
the parabolic Lie algebroid $(V_{*},\, \phi_{*})$. It should be clarified that we can choose the
ramified Galois covering map $\sigma$ in such a way that any given finite set of parabolic bundles on
$(X,\,S)$ correspond to equivariant bundles on $Y$. To see this, take ramified Galois coverings for
the individual parabolic  bundles in the given finite set of parabolic
bundles. Take the fiber product of
these finitely many ramified Galois coverings, and then take its normalization. The resulting
ramified Galois covering satisfies the condition that each of the parabolic  bundles in the
given finite set corresponds to an equivariant  bundle on the covering. This shows that $\mathcal{E}_{G}$ and $(V_{*},\phi_{*})$ can simultaneously be realized as an equivariant principal $G-$bundle and an equivariant holomorphic Lie algebroid respectively.

We note that $\phi$ is not surjective because $\phi_{*}$ is not surjective.

By Theorem \ref{4.1} we know that there is a $\Gamma$--equivariant holomorphic $(V,\,\phi)$--connection on 
$E_{P}$. Therefore, applying Theorem \ref{char}(v) for the complex linear algebraic group $P$, it follows 
that there is a parabolic $(V_{*},\, \phi_{*})$--connection on $\mathcal{E}_{P}$.
\end{proof}

\end{document}
